\chapter{結果}
\label{cha:results}

\section{表面下散乱分布の計測結果}
\label{sec:results_rd}

本節では，
前章で確立した評価手法を用いて得られた
自作レンダラーによる計測結果を示す。
\begin{figure}[H]
  \centering
  \includegraphics[width=0.8\linewidth]{figs/chot.pdf}
  \caption{髭なし，髭モデル，混合モデルにおける正規化表面下散乱分布 $R_d(r)$ の比較。}
  \label{fig:rd_norm_compare}
\end{figure}

図\ref{fig:rd_norm_compare}は，
髭なしの皮膚モデル，
髭モデル，
および両者を混合したモデルに対して計測した
正規化後の表面下散乱分布 $R_d(r)$ を比較したものである。
いずれの分布も面積積分が 1 となるよう正規化されており，
分布形状の違いに着目した比較が可能である。

図より，
各条件において $R_d(r)$ の分布形状に差が見られ，
特に混合モデルの分布は，
髭なしモデルと髭モデルの中間的な挙動を示している。
また，
いずれの分布も距離 $r$ に対して滑らかに変化しており，
計測結果がノイズや発散に支配されていないことが確認できる。

これらの結果は，
髭の有無やモデルの違いが
表面下散乱分布に反映されていることを示している。
次節では，
これらの分布の違いについて考察を行う。
\section{考察}

\section{幾何的な髭モデルによる検証の限界}

本研究では当初，
髭を円柱などの幾何形状として明示的に配置したシーンを用い，
表面下散乱の放射輝度分布を計測することで，
髭が皮膚表面下散乱に与える影響を直接的に評価することを試みた。

しかし，
このような幾何的髭モデルによる検証では，
いくつかの課題が明らかとなった。

まず，
髭の配置位置や深さ，太さのわずかな違いによって，
計測される放射輝度分布が大きく変動する傾向が確認された。
これは，
髭が作る自己遮蔽や影，
さらに皮膚内部での多重散乱との相互作用が複雑に絡み合うためであり，
幾何的に厳密なモデルであっても，
安定した再現性を得ることが困難であることを示している。

また，
髭一本一本を幾何形状として扱う場合，
シーン内のプリミティブ数が急激に増加し，
計算コストが大きくなるという実用上の問題も無視できない。
このため，
幾何モデルは理論的には詳細な表現が可能である一方，
実際のレンダリングや検証においては，
不安定かつ高コストな手法となる。

以上より，
幾何的な髭モデルは，
「髭が表面下散乱に影響を与えうる」ことを確認する手段としては有効であるが，
その影響を安定的かつ制御可能な形で評価する基準モデルとして用いることは，
必ずしも適切ではないと考えられる。

\section{髭密度モデルによる表面下散乱制御の妥当性}

これに対し，
本研究で提案した髭密度モデルは，
髭の幾何構造そのものを直接扱うのではなく，
髭の分布を表面上のスカラー量として近似し，
皮膚モデルと髭モデルを線形に混合するという設計を採用している。

第\ref{sec:results_rd}節で示した結果より，
この髭密度に基づく混合モデルでは，

\begin{itemize}
  \item 放射輝度分布 $R_d(r)$ が不自然に破綻することなく，
  \item 髭密度の増加に応じて，
        表面下散乱の広がりが連続的に変化する
\end{itemize}

ことが確認された。

特に注目すべき点は，
放射輝度分布を正規化した後においても，
髭密度の違いによる分布形状の差が明確に現れている点である。
これは，
単に全体のエネルギー量が変化しているのではなく，
光の空間的な拡散特性そのものが変化していることを示している。

この結果は，
髭密度という単純化された指標を用いた線形混合であっても，
表面下散乱の挙動を十分に特徴づけられることを示唆しており，
提案モデルの設計判断が妥当であることを支持している。

\section{幾何モデルと密度モデルの関係性について}

本研究において重要なのは，
髭密度モデルが幾何的髭モデルを厳密に再現することではない。
むしろ，
髭が表面下散乱に与える影響は，

\begin{itemize}
  \item 髭の太さ
  \item 配置密度
  \item 皮膚内部への侵入深さ
  \item 髭自身が作る影や遮蔽
\end{itemize}

といった多くの要因が複雑に関係しており，
これらをすべて厳密にモデル化することは現実的ではない。

このような背景を踏まえると，
髭密度モデルは，
幾何的に複雑な要因を一つのスカラー量に集約し，
その影響を現象論的に捉えるための抽象化であると位置付けられる。

第\ref{sec:results_rd}節で示した差分実験の結果から，
幾何的髭モデルと髭密度モデルの間には，
放射輝度分布の変化傾向において共通点が見られた。
このことは，
密度モデルが幾何モデルの詳細を再現していないにもかかわらず，
髭の存在による影響を方向性として適切に捉えていることを示している。

したがって，
提案モデルは，
「物理的に完全な髭モデル」ではなく，
「表面下散乱に対する髭の影響を安定かつ制御可能な形で表現するための実用的モデル」
として位置付けることができる。

\section{本研究の位置付けと限界}

本研究で提案した髭密度モデルは，
計算コストと表現力のバランスを重視した近似モデルであり，
すべての髭表現を正確に再現できるものではない。
特に，
個々の髭が作る鋭い影や，
髭同士の相互遮蔽といった効果については，
本モデルでは明示的に扱っていない。

一方で，
顔全体のレンダリングやインタラクティブな応用を想定した場合，
髭の影響を密度という形で扱い，
表面下散乱の広がりを連続的に制御できる点は，
実用上大きな利点である。

以上より，
本研究のモデルは，
高精度な物理再現を目的とするモデルとは異なる立場にあり，
実用的な皮膚レンダリングにおける
新たな設計指針を提示するものであると考えられる。

\section{まとめ}

本章では，
幾何的な髭モデルと比較しながら，
髭密度に基づく表面下散乱モデルの妥当性について考察を行った。

その結果，

\begin{itemize}
  \item 髭は表面下散乱の放射輝度分布に影響を与えること
  \item 幾何モデルによる厳密な評価は不安定かつ高コストであること
  \item 髭密度モデルは，
        その影響を安定的かつ連続的に制御可能であること
\end{itemize}

が確認された。

これらの知見は，
髭のような微細構造を持つ要素を含む皮膚表現において，
密度に基づく近似モデルが有効な選択肢となり得ることを示している。


